\documentclass[border=8pt]{standalone}
\usepackage[RPvoltages]{circuitikz}

% Propuesta T3, segundo semestre de 2026.
% Convenciones tomadas de electronic_circuits_diagrams/src.
% Compilar con: pdflatex T3_howland_invertida.tex
% La resistencia de seguridad forma parte del circuito desde el inicio.

\begin{document}
\begin{circuitikz}
\ctikzset{resistors/scale=0.8}
\draw (0,0) node[op amp, noinv input up](OA){};

% Nodo de carga L: tambien es la entrada no inversora del operacional.
\draw (-5.5,2) node[ground]{}
    to[R, l=$R_{19}$, -*] (-3,2) coordinate(L)
    to[R, l=$R_{20}$] (3,2) -- (3,0) coordinate(A);
\draw (L) -- (-3,0.5) |- (OA.+);
\node[below left] at (L) {$v_{L,N}$};

% Divisor conectado a la entrada inversora.
% En la fuente invertida se intercambian vi y tierra, NO los pines + y -.
\draw (-5.5,-2) node[below]{$v_i$}
    to[R, l_=$R_{21}$, o-*] (-3,-2) coordinate(M)
    to[R, l_=$R_{22}$] (3,-2)
    to[short, -*] (A);
\draw (M) -- (-3,-0.5) |- (OA.-);
\draw (OA.out) to[short, -*] (A) -- (3.8,0)
    node[right]{$v_{A,N}$};

% Alimentacion del operacional; tierra es 0 V.
\draw (OA.up) -- ++(0,0.15) node[vcc]{$+2.5\,\mathrm{V}$};
\draw (OA.down) -- ++(0,-0.15) node[vee]{$-2.5\,\mathrm{V}$};

% Carga: anodo arriba, catodo al riel negativo.
% La carga se toma de L, NO de la salida A del operacional.
\draw (L) -- (-3,3.6) -- (5.5,3.6);
\draw (5.5,3.6) to[R, l=$R_{23}$] (5.5,1.4);
\draw (5.5,1.4) to[short, i=$i_N$] (5.5,0.8)
    to[leDo, l=$\mathrm{LED7}$] (5.5,-1.2)
    -- (5.5,-2.3) node[vee]{$-2.5\,\mathrm{V}$};
    
\end{circuitikz}
\end{document}
